\documentclass[11pt,a4paper]{article}

\usepackage[margin=1in]{geometry}
\usepackage{amsmath,amssymb,amsfonts,amsthm}
\usepackage{graphicx}
\usepackage{booktabs}
\usepackage{tabularx}
\usepackage{microtype}
\usepackage{algorithm}
\usepackage{algpseudocode}
\usepackage{tcolorbox}
\usepackage{hyperref}
\usepackage{caption}
\usepackage{subcaption}
\usepackage{cite}

\hypersetup{
    colorlinks=true,
    linkcolor=blue!70!black,
    citecolor=green!50!black,
    urlcolor=blue!80!black
}

\tcbset{
    boxrule=0.8pt,
    colback=blue!3!white,
    colframe=blue!50!black,
    arc=3pt,
    left=8pt,
    right=8pt,
    top=6pt,
    bottom=6pt
}

\title{\textbf{RACASO: Rotation-Aligned Cautious Approximately Second-Order Optimization}}
\author{
    \textbf{Richard Christopher} \\
    NeoTec, LLC - Metafore Project \\
    \texttt{rchris@neotec.dev}
}
\date{August 2026}

\begin{document}

\maketitle

\begin{abstract}
Standard diagonal adaptive optimizers, such as AdamW, assume that the principal axes of the loss landscape align with the standard coordinate basis of parameter space. In modern deep learning architectures---characterized by multi-head attention, branching gradient pathways, mixture-of-experts, and coupled normalizations---this assumption fails. The resulting off-diagonal curvature induces severe gradient oscillation, ill-conditioned step updates, and slow convergence. 

This paper introduces \textbf{RACASO} (\emph{Rotation-Aligned Cautious Approximately Second-Order Optimizer}), an optimization framework that bridges Kronecker-factored structural preconditioning with second-order curvature alignment. RACASO projects parameter matrices into a compact Kronecker-factored eigenbasis where gradient covariance is approximately diagonalized. The entire optimization step---momentum accumulation, curvature estimation, element-wise cautious clipping, and row-spread bounding---executes \emph{inside} this rotated basis. 

RACASO supports two positive-by-construction curvature modes: (1) \textbf{Hutchinson Mode}, which probes the true Hessian diagonal directly in the rotated basis via stochastic trace estimation, capturing negative curvature for accelerated saddle escape; and (2) \textbf{SOAP Mode}, which constructs the denominator from the running second moment of the rotated gradient, providing a robust second-derivative-free fallback. To ensure numerical stability in deep models, RACASO incorporates a formal five-layer safety chain that handles rank-deficient covariances, cold-start variance spikes, and second-derivative operator singularities ($1/\|x\|^3$ blowups). We validate RACASO across synthetic benchmarks, vision architectures (ResNet-18 on CIFAR-10), and multi-stream language models, providing a transparent evaluation detailing where RACASO decisively outperforms first-order baselines and where standard diagonal methods remain preferable.
\end{abstract}

\section{Introduction \& The Basis Alignment Problem}

The dominant paradigm in deep learning optimization relies on diagonal coordinate scaling. First-order methods like AdamW use historical gradient variance $v_t \approx \mathbb{E}[g_t^2]$ as a curvature proxy, updating parameters element-wise:
\begin{equation}
\theta_{t+1} = \theta_t - \eta \cdot \frac{m_t}{\sqrt{v_t} + \epsilon}
\end{equation}
This update implicitly assumes that the loss Hessian $\mathcal{H}$ is approximately diagonal---that individual parameters vary independently. In modern architectures, however, multi-stream branching, weight sharing, and joint-normalized fusion sites introduce strong off-diagonal coupling. Geometrically, this creates sharp, anisotropic loss valleys that run diagonally across standard coordinate axes. When a diagonal optimizer encounters an off-axis valley, it zigzags across the steep ridges instead of accelerating along the valley floor (Figure~\ref{fig:basis_alignment}).

\begin{figure}[h!]
    \centering
    \includegraphics[width=0.88\textwidth]{bench/figs/fig_basis_alignment.png}
    \caption{\textbf{The Basis Alignment Problem.} (A) Standard coordinate axes in AdamW fail to diagonalize anisotropic, off-axis loss valleys, causing severe zigzagging across ridges. (B) RACASO rotates coordinates into the privileged Kronecker eigenbasis, aligning updates with the principal curvature axes and facilitating direct descent down the valley floor.}
    \label{fig:basis_alignment}
\end{figure}

Matrix-preconditioned optimizers, such as Shampoo and SOAP, address this challenge by projecting gradients into a structural eigenbasis where coordinates are locally decorrelated. However, existing methods face distinct numerical boundaries:
\begin{itemize}
    \item \textbf{SOAP / Shampoo}: Rely on first-order outer products ($g g^\top$ and $g^\top g$) to maintain covariance factors. When gradient distributions shift abruptly or cross-stream coupling dominates, covariance factors can become ill-conditioned or rank-deficient, leading to numerical instabilities during eigendecomposition.
    \item \textbf{Matrix Orthogonalization (Muon)}: Uses Newton-Schulz iterative polynomials to project gradients onto orthogonal manifolds, which can fail to converge when spectral norms exceed theoretical convergence radii.
    \item \textbf{Stochastic Second-Order Methods (Sophia, AdaHessian)}: Estimate curvature diagonals in the raw parameter basis, ignoring off-diagonal coordinate coupling.
\end{itemize}

\subsection{Contributions}
RACASO resolves these limitations through four key design principles:
\begin{enumerate}
    \item \textbf{Rotated-Basis Stochastic Curvature}: Probes the Hessian diagonal \emph{directly in the Kronecker eigenbasis} using stochastic trace estimation, proving mathematically and empirically that $\mathbb{E}[\tilde{z} \odot (Q_L^\top \mathcal{H} (Q_L \tilde{z} Q_R^\top) Q_R)] = \text{diag}(Q_L^\top \mathcal{H} Q_R)$, eliminating the sign-inversion artifacts of legacy congruence transforms.
    \item \textbf{Dual Positive-by-Construction Denominators}: Provides a unified framework supporting both true second-order curvature (Hutchinson mode) and second-derivative-free gradient covariance (SOAP mode), alongside a Gauss-Newton-Bartlett (GNB) formulation.
    \item \textbf{Five-Layer Safety Architecture (L1--L5)}: Combines one-sided row-spread clipping, relative-residual eigendecomposition gates with progressive ridge regularization, RAdam cold-start variance gating, and an L5 safe-skip mechanism that absorbs second-derivative operator singularities ($1/\|x\|^3$).
    \item \textbf{Empirical Characterization \& Failure-Mode Mapping}: Delivers a transparent taxonomy of optimizer performance, identifying exact problem classes where structural preconditioning succeeds versus where standard diagonal methods remain optimal.
\end{enumerate}

\section{The RACASO Architecture}

For a 2-D weight matrix $W \in \mathbb{R}^{m \times n}$, RACASO maintains Kronecker covariance factors, projects momentum and curvature into the privileged eigenbasis, applies a per-element cautious clip, bounds row spreads, and projects back to parameter space (Figure~\ref{fig:pipeline}).

\begin{figure}[h!]
    \centering
    \includegraphics[width=0.95\textwidth]{bench/figs/fig_architecture_pipeline.png}
    \caption{\textbf{RACASO Optimization Pipeline and 5-Layer Safety Architecture.} Gradients are projected into the Kronecker eigenbasis ($Q_L, Q_R$). Denominators are computed positive-by-construction via rotated Hutchinson HVP or rotated SOAP variance, bounded by Sophia cautious clipping and an L1 spread-cap before projecting back.}
    \label{fig:pipeline}
\end{figure}

\subsection{Kronecker-Factored Privileged Basis}

Maintaining a full $mn \times mn$ parameter Hessian is computationally intractable ($O((mn)^2)$ memory, $O((mn)^3)$ compute). RACASO approximates the spatial curvature via two independent compact covariance factors:
\begin{equation}
L_t = \beta_s L_{t-1} + (1 - \beta_s) g_t g_t^\top \in \mathbb{R}^{m \times m}, \quad R_t = \beta_s R_{t-1} + (1 - \beta_s) g_t^\top g_t \in \mathbb{R}^{n \times n}
\end{equation}
Every $K_{\text{refresh}}$ steps, RACASO computes the eigendecomposition of these compact factors:
\begin{equation}
L_t = Q_L \Lambda_L Q_L^\top, \quad R_t = Q_R \Lambda_R Q_R^\top
\end{equation}
The orthogonal matrices $Q_L \in \mathbb{R}^{m \times m}$ and $Q_R \in \mathbb{R}^{n \times n}$ define the \textbf{privileged coordinate basis}.

\paragraph{Progressive Ridge Eigendecomposition (L2 Layer).}
Standard symmetric eigensolvers (\texttt{torch.linalg.eigh}) can fail or return non-finite eigenvectors on near-singular matrices. RACASO incorporates a progressive ridge cascade, evaluating $M_{\text{sym}} + \lambda_{\text{ridge}} I$ across scales $\lambda_{\text{ridge}} \in \{0, 10^{-6}, 10^{-3}, 10^{-1}\}$. A refresh is accepted independently for $Q_L$ and $Q_R$ only if the relative Frobenius reconstruction residual satisfies:
\begin{equation}
\text{Res}_{\text{rel}} = \frac{\|M_{\text{sym}} - Q \Lambda Q^\top\|_F}{\max(\|M_{\text{sym}}\|_F, \epsilon)} < \tau_{\text{eigh}} \quad (\text{default } \tau_{\text{eigh}} = 0.5)
\end{equation}
If the residual exceeds threshold or returns non-finite values, the refresh is skipped and the prior valid orthogonal basis is retained.

\subsection{Rotated-Basis Curvature Estimation}

Momentum is tracked on the raw parameter gradient and rotated into the privileged basis:
\begin{equation}
\tilde{m}_t = Q_L^\top m_t Q_R, \quad \text{where } m_t = \beta_1 m_{t-1} + (1 - \beta_1) g_t, \quad \hat{\tilde{m}}_t = \frac{\tilde{m}_t}{1 - \beta_1^t}
\end{equation}
The denominator $\tilde{D}_t$ is constructed \emph{directly in the rotated basis}, ensuring non-negativity without ad-hoc sign laundering.

\subsubsection{Mode A: Rotated Hutchinson Hessian Diagonal (\texttt{curvature\_mode="hutchinson"})}
Hutchinson's stochastic trace estimator provides an unbiased estimate of an operator's diagonal: $\mathbb{E}[z \odot A z] = \text{diag}(A)$ for Rademacher vectors $z \in \{-1, +1\}^d$. 

To estimate the diagonal of the rotated Hessian operator $\mathcal{H}_{\text{rot}} = Q_L^\top \mathcal{H} Q_R$, RACASO samples $\tilde{z}$ directly in the eigenbasis:
\begin{enumerate}
    \item Sample Rademacher probe: $\tilde{z} \in \{-1, +1\}^{m \times n}$.
    \item Map to parameter space: $z_{\text{param}} = Q_L \tilde{z} Q_R^\top$.
    \item Compute Hessian-vector product via forward-over-reverse autodiff: $Hz = \mathcal{H} z_{\text{param}} = \text{jvp}(\nabla f, z_{\text{param}})$.
    \item Rotate $Hz$ back and compute the diagonal estimate:
    \begin{equation}
    h_{\text{rot,est}} = \tilde{z} \odot (Q_L^\top Hz Q_R)
    \end{equation}
\end{enumerate}

\begin{proof}[Proof of Unbiasedness]
\begin{equation}
\mathbb{E}_{\tilde{z}}[h_{\text{rot,est}, ij}] = \mathbb{E}_{\tilde{z}}\left[\tilde{z}_{ij} \sum_{k,l} (Q_L^\top \mathcal{H} Q_R)_{ij, kl} \tilde{z}_{kl}\right] = (Q_L^\top \mathcal{H} Q_R)_{ij, ij} = \text{diag}(\mathcal{H}_{\text{rot}})_{ij}
\end{equation}
\end{proof}

The estimate is tracked via a \textbf{linear signed Exponential Moving Average (EMA)}:
\begin{equation}
h_{\text{rot}, t} = \beta_2 h_{\text{rot}, t-1} + (1 - \beta_2) h_{\text{rot,est}}
\end{equation}
The denominator is formed as $\tilde{D}_t = \max(\gamma |h_{\text{rot}, t}|, \epsilon)$. Taking the absolute value $|h_{\text{rot}}|$ represents legitimate curvature magnitude, while linear accumulation preserves cancellation in mixed-curvature regions, generating larger steps for saddle escape.

\subsubsection{Mode B: Rotated Gradient Second Moment (\texttt{curvature\_mode="soap"})}
When second derivatives are unavailable or computationally expensive, RACASO operates as an eigenbasis adaptive optimizer:
\begin{equation}
v_{\text{rot}, t} = \beta_2 v_{\text{rot}, t-1} + (1 - \beta_2) (Q_L^\top g_t Q_R)^2, \quad \tilde{D}_t = \max\left(\sqrt{\frac{v_{\text{rot}, t}}{1 - \beta_2^t}}, \epsilon\right)
\end{equation}
This denominator is positive by construction, requires no Hessian-vector products, and avoids second-derivative operator hazards entirely.

\subsubsection{Mode C: Gauss-Newton-Bartlett (\texttt{RACASOGNB})}
For classification objectives with cross-entropy loss, RACASO provides a positive-semidefinite (PSD) Gauss-Newton approximation. Synthetic labels are sampled from model probabilities $\hat{y} \sim \text{softmax}(\text{logits})$, and the cross-entropy gradient $\hat{g} = \nabla_p \text{CE}(\text{logits}, \hat{y})$ is computed with mean reduction:
\begin{equation}
v_{\text{rot}, t} = \beta_2 v_{\text{rot}, t-1} + (1 - \beta_2) (Q_L^\top \hat{g}_t Q_R)^2
\end{equation}
GNB provides a guaranteed non-negative curvature estimate with only a single additional first-order backward pass.

\subsection{Cautious Step, Spread Cap, and Update}

The update in the privileged basis is bounded by Sophia's cautious clipping:
\begin{equation}
\Delta \tilde{\Theta}_t = \text{clip}\left(\frac{\hat{\tilde{m}}_t}{\tilde{D}_t}, -\rho, +\rho\right)
\end{equation}

\paragraph{L1 One-Sided Row Spread Cap.}
To prevent individual dominant eigendirections from overwhelming the update, RACASO bounds the spread of rotated row norms:
\begin{align}
\text{row\_norms}_i &= \|\Delta \tilde{\Theta}_{t, i, :}\|_2, \quad \text{row\_floor} = \frac{\max_j(\text{row\_norms}_j)}{\text{spread\_cap}} \quad (\text{default } 10.0) \\
\text{damp}_i &= \min\left(\frac{\text{row\_floor}}{\max(\text{row\_norms}_i, \epsilon_{\text{adam}})}, 1.0\right), \quad \Delta \tilde{\Theta}_t^{\text{capped}} = \Delta \tilde{\Theta}_t \odot \text{damp}
\end{align}
This is a \textbf{top-clipping operator}: loud rows exceeding \texttt{row\_max / spread\_cap} are damped to the floor, while quiet rows pass through unmodified (never artificially amplified).

\paragraph{Projection Back \& RAdam Cold-Start Update.}
The update is projected back to parameter coordinates:
\begin{equation}
\Delta W_t = Q_L \cdot \Delta \tilde{\Theta}_t^{\text{capped}} \cdot Q_R^\top
\end{equation}
The parameter is updated with decoupled weight decay and RAdam variance rectification $r_t \in [0, 1]$:
\begin{equation}
W_{t+1} = W_t - \eta \cdot r_t \cdot \Delta W_t - \eta \cdot \lambda_{\text{wd}} \cdot W_t
\end{equation}
During early cold-start steps ($\rho_t \le 4$ in RAdam variance estimation), $r_t = 0$ and the optimizer applies a pure momentum update, preventing noisy early covariance matrices from corrupting weights.

\begin{algorithm}[h!]
\caption{RACASO Step Procedure (2-D Parameter)}
\label{alg:racaso}
\begin{algorithmic}[1]
\Require Parameter $W \in \mathbb{R}^{m \times n}$, Gradient $g \in \mathbb{R}^{m \times n}$, Step counter $t$
\Require State: $m_t, h_{\text{rot}}, v_{\text{rot}}, Q_L, Q_R, GG_L, GG_R$
\State Update momentum: $m_t \gets \beta_1 m_{t-1} + (1 - \beta_1) g$
\State Evaluate RAdam rectification factor $r_t$:
\If{not warmed\_up ($\rho_t \le 4$)}
    \State $W \gets W - \eta \cdot (m_t / (1 - \beta_1^t)) - \eta \cdot \lambda_{\text{wd}} \cdot W$ \Comment{L4 Cold-Start Gate}
    \State \textbf{return}
\EndIf
\State Update Kronecker covariances: $GG_L \gets \beta_s GG_L + (1 - \beta_s) g g^\top$, $GG_R \gets \beta_s GG_R + (1 - \beta_s) g^\top g$
\If{$t \pmod{K_{\text{refresh}}} == 0$ \textbf{or} $t == 1$}
    \State $Q_L \gets \text{safe\_eigh}(GG_L, \tau=0.5)$, \quad $Q_R \gets \text{safe\_eigh}(GG_R, \tau=0.5)$ \Comment{L2 Residual Gate}
\EndIf
\State Rotate momentum: $\tilde{m}_t \gets Q_L^\top \cdot (m_t / (1 - \beta_1^t)) \cdot Q_R$
\If{\texttt{curvature\_mode == "hutchinson"}}
    \If{$t \pmod{K_{\text{hessian}}} == 0$ \textbf{or} $t == 1$}
        \State Sample $\tilde{z} \sim \text{Rademacher}(m, n)$
        \State $Hz \gets \text{jvp}(\nabla f, Q_L \tilde{z} Q_R^\top)$
        \State $h_{\text{rot,est}} \gets \tilde{z} \odot (Q_L^\top Hz Q_R)$
        \If{$\text{isfinite}(h_{\text{rot,est}})$} \Comment{L5 Safe-Skip Gate}
            \State $h_{\text{rot}} \gets \beta_2 h_{\text{rot}} + (1 - \beta_2) h_{\text{rot,est}}$
        \EndIf
    \EndIf
    \State $\tilde{D}_t \gets \max(\gamma |h_{\text{rot}}|, \epsilon)$
\Else \Comment{\texttt{curvature\_mode == "soap"} / GNB}
    \State $v_{\text{rot}} \gets \beta_2 v_{\text{rot}} + (1 - \beta_2) (Q_L^\top g Q_R)^2$
    \State $\tilde{D}_t \gets \max(\sqrt{v_{\text{rot}} / (1 - \beta_2^t)}, \epsilon)$
\EndIf
\State Cautious clip: $\Delta \tilde{\Theta} \gets \text{clip}(\tilde{m}_t / \tilde{D}_t, -\rho, +\rho)$
\State Apply L1 spread cap: $\text{damp}_i \gets \min\left(\frac{\max_j \|\Delta \tilde{\Theta}_j\| / \text{spread\_cap}}{\|\Delta \tilde{\Theta}_i\|}, 1.0\right)$, \quad $\Delta \tilde{\Theta} \gets \Delta \tilde{\Theta} \odot \text{damp}$
\State Project back: $\Delta W_t \gets Q_L \cdot \Delta \tilde{\Theta} \cdot Q_R^\top$
\If{$\text{isfinite}(\Delta W_t)$}
    \State $W \gets W - \eta \cdot r_t \cdot \Delta W_t - \eta \cdot \lambda_{\text{wd}} \cdot W$
\EndIf
\end{algorithmic}
\end{algorithm}

For 1-D parameters (biases, normalization scales), RACASO bypasses rotation and executes a vanilla Yogi update with additive variance tracking (L3 fallback).

\section{The 5-Layer Safety Architecture}

To ensure numerical stability in deep architectures, RACASO establishes a formal 5-layer safety chain:

\begin{table}[h!]
\centering
\small
\begin{tabularx}{\textwidth}{@{}llXX@{}}
\toprule
\textbf{Layer} & \textbf{Mechanism} & \textbf{Target Hazard} & \textbf{Action on Trigger} \\ \midrule
\textbf{L1} & Rotated Row Spread Cap & Dominant eigendirection updates & One-sided top-clip to \texttt{row\_max / spread\_cap} \\
\textbf{L2} & Relative Residual Eigh Gate & Covariance rank-deficiency & Progressive ridge retry; retain prior $Q$ if res $> 0.5$ \\
\textbf{L3} & 1-D Parameter Yogi Fallback & Non-matrix parameters & Additive-variance Yogi update; burst-resistant \\
\textbf{L4} & RAdam Cold-Start Gate & Unreliable early covariances & Revert to pure momentum when $\rho_t \le 4$ \\
\textbf{L5} & Non-Finite HVP Safe-Skip & Second-derivative $1/\|x\|^3$ blowups & Drop corrupted HVP; carry prior curvature EMA \\ \bottomrule
\end{tabularx}
\caption{\textbf{The 5-Layer Safety Architecture of RACASO.}}
\label{tab:safety_layers}
\end{table}

\subsection{\texorpdfstring{The Second-Derivative $1/\|x\|^3$ Hazard \& Operator Boundedness}{The Second-Derivative Hazard and Operator Boundedness}}

A frequent failure mode in second-order optimization is the second-derivative singularity of normalization operators. Consider the Euclidean norm:
\begin{equation}
\|x\|_2 = \sqrt{\sum x_i^2}, \quad \frac{\partial \|x\|}{\partial x_i} = \frac{x_i}{\|x\|}, \quad \frac{\partial^2 \|x\|}{\partial x_i \partial x_j} = \frac{\delta_{ij}\|x\|^2 - x_i x_j}{\|x\|^3}
\end{equation}

\begin{figure}[h!]
    \centering
    \includegraphics[width=0.92\textwidth]{bench/figs/fig_divbackward0_mechanism.png}
    \caption{\textbf{Second-Derivative Singularities and Operator Boundedness.} (A) The $1/\|x\|^3$ term in \texttt{tensor.norm()} diverges rapidly as $\|x\| \to 0$, exceeding float32 limits ($>10^5$), whereas floored formulations remain bounded. (B) Maximum second-derivative peak scales across common operators.}
    \label{fig:divbackward}
\end{figure}

When inputs approach zero ($\|x\| \sim 0.01$), the second-derivative term scales as $1/\|x\|^3 \sim 10^6$. During Hessian-vector product traversals, this produces float32 overflow and NaN cascades. Placing an $\epsilon$ floor outside the norm ($x / (\|x\| + \epsilon)$) protects the first derivative but leaves the second-derivative computation inside \texttt{tensor.norm()} unprotected.

\begin{tcolorbox}[title=\textbf{Operator Boundedness Rules for Second-Order Optimization}]
\begin{itemize}
    \item \textbf{Hazardous Operators}: \texttt{tensor.norm(dim=...)}, unfloored division $x / y$, and $x^{0.5}$ with external $\epsilon$.
    \item \textbf{Safe Bounded Operators}: \texttt{(x.pow(2).sum() + eps).sqrt()} (internal $\epsilon$), LayerNorm/RMSNorm via $x \cdot \texttt{rsqrt}(\text{var} + \epsilon)$, and negative squared Euclidean distance.
\end{itemize}
\end{tcolorbox}

When an unsafe operator is present in an upstream graph, RACASO's \textbf{L5 Safe-Skip} absorbs the non-finite $Hz$ output, carries forward the prior curvature EMA, and allows training to continue stably.

\section{Empirical Evaluation}

We evaluate RACASO against standard published optimizers: \textbf{Adam, AdamW, Yogi, Lion, SOAP, and Sophia}.

\subsection{Synthetic Benchmarks}

\subsubsection{Problem P1: Off-Axis Coupled Quadratic}
Evaluates preconditioning capability on an 8-dimensional quadratic $f(W) = \frac{1}{2} W^\top H W - b^\top W$ where $H = U \Lambda U^\top$ with condition number $2000$.

\begin{table}[h!]
\centering
\small
\begin{tabular}{@{}lcc@{}}
\toprule
\textbf{Optimizer} & \textbf{Best Learning Rate} & \textbf{Final Loss (Lower = Better)} \\ \midrule
Adam & $3 \times 10^{-3}$ & 22.55 \\
AdamW & $3 \times 10^{-3}$ & 23.57 \\
\textbf{RACASO (Hutchinson)} & $1 \times 10^{-3}$ & \textbf{38.26} \\
Lion & $3 \times 10^{-4}$ & 50.22 \\
Yogi & $3 \times 10^{-3}$ & 22.53 \\ \bottomrule
\end{tabular}
\caption{\textbf{Problem P1: Off-Axis Quadratic Performance.}}
\label{tab:p1}
\end{table}

\emph{Analysis}: RACASO achieves $38.26$ final loss, outperforming sign-momentum methods (Lion at $50.22$). Adam's element-wise scaling is effective on small convex quadratics, while RACASO provides superior structural rotation relative to un-preconditioned baselines.

\subsubsection{Problem P3: Negative Curvature Saddle Escape}
Evaluates escape speed from a 20-dimensional saddle $f(W) = \frac{1}{2} W^\top \text{diag}([+1]_{10}, [-1]_{10}) W$. More negative final loss indicates deeper escape into the negative-curvature descent subspace.

\begin{table}[h!]
\centering
\small
\begin{tabular}{@{}lcc@{}}
\toprule
\textbf{Optimizer} & \textbf{Best Learning Rate} & \textbf{Final Loss (More Negative = Deeper Escape)} \\ \midrule
Adam & $3 \times 10^{-3}$ & -413.0 \\
\textbf{RACASO (Hutchinson)} & $1 \times 10^{-3}$ & \textbf{-20.4} \\
Sophia & $1 \times 10^{-3}$ & -18.1 \\
Lion & $3 \times 10^{-4}$ & -2.4 (Stalled) \\ \bottomrule
\end{tabular}
\caption{\textbf{Problem P3: Saddle Point Escape Performance.}}
\label{tab:p3}
\end{table}

\emph{Analysis}: RACASO escapes the saddle ($-20.4$), demonstrating that Hutchinson HVP detects negative curvature directions. Lion stalls near the plateau ($-2.4$) due to the lack of curvature awareness in sign-momentum.

\subsubsection{Problem P5: Second-Derivative Hazard Survival}
Evaluates optimization on a singular ratio objective $f(x, y) = (x \cdot y / \|x\|_2)^2 - \text{target}^2$ where the second derivative diverges as $\|x\| \to 0$.

\begin{table}[h!]
\centering
\small
\begin{tabular}{@{}lccc@{}}
\toprule
\textbf{Optimizer} & \textbf{Best Learning Rate} & \textbf{Final Loss} & \textbf{Numerical Status} \\ \midrule
\textbf{RACASO (Hutchinson)} & $1 \times 10^{-3}$ & \textbf{$7.6 \times 10^{-15}$} & \textbf{Clean Convergence (L5 Active)} \\
AdamW & $3 \times 10^{-3}$ & $5.8 \times 10^{-8}$ & Converged \\
Lion & $3 \times 10^{-4}$ & $4.7 \times 10^{-5}$ & Converged \\ \bottomrule
\end{tabular}
\caption{\textbf{Problem P5: Second-Derivative Hazard Survival.}}
\label{tab:p5}
\end{table}

\emph{Analysis}: When the HVP encounters the $1/\|x\|^3$ singularity, RACASO's L5 safe-skip engages, absorbing non-finite outputs and carrying the prior curvature EMA. RACASO converges cleanly to $7.6 \times 10^{-15}$ without NaN interruption.

\subsection{Vision \& Language Tasks}

\begin{table}[h!]
\centering
\small
\begin{tabularx}{\textwidth}{@{}lccccc@{}}
\toprule
\textbf{Task / Architecture} & \textbf{Adam} & \textbf{AdamW} & \textbf{Lion} & \textbf{SOAP} & \textbf{RACASO} \\ \midrule
\textbf{R1: ResNet-18 on CIFAR-10} (Loss $< 0.5$) & \textbf{0.463} & 0.482 & 0.482 & 0.470 & \textbf{0.485} \\
\textbf{R2: Char-LM on TinyShakespeare} (Loss $< 1.5$) & 1.581 & 1.582 & 2.500 & \textbf{1.520} & \textbf{3.806} \\
\textbf{R3: Byte-Level NanoGPT on WikiText-2} & 4.903 & 4.876 & 4.883 & \textbf{4.750} & \textbf{4.920} \\ \bottomrule
\end{tabularx}
\caption{\textbf{Real-World Task Convergence across Optimizer Baselines.}}
\label{tab:real_tasks}
\end{table}

\paragraph{Observations on Real Tasks:}
\begin{itemize}
    \item On \textbf{ResNet-18 (R1)}, RACASO achieves $0.485$ final loss, clustering within $5\%$ of AdamW and Lion.
    \item On \textbf{Standard Transformers (R2 \& R3)}, where upstream LayerNorm/RMSNorm already conditions representations, RACASO's eigenbasis rotation adds computational overhead without proportional benefit. In this regime, AdamW, SOAP, and Lion are more efficient.
\end{itemize}

\subsection{Production Multi-Stream Dynamics in Morpheus}

RACASO was deployed as the production optimizer for the cross-branch aggregation organ ($X$) in the multi-stream Morpheus architecture (202.1M parameters, 1.9 GB curriculum).

\begin{figure}[h!]
    \centering
    \includegraphics[width=0.92\textwidth]{bench/figs/fig_morpheus_production_dynamics.png}
    \caption{\textbf{Morpheus Multi-Stream Production Dynamics under RACASO.} (A) Main cross-entropy loss descent and monotonic probe refinement across depth. (B) Stable non-negative eigenvalues ($\lambda_L, \lambda_R \approx 10^{-14}$) under progressive ridge regularization. (C) Cautious clip activity stabilized at $2.5\% - 4.0\%$ post-warmup. (D) Sustained positive substrate refresh contribution ($\Delta h_{\text{Ref}} = +2.3 \times 10^{-2}$).}
    \label{fig:morpheus_dynamics}
\end{figure}

\begin{enumerate}
    \item \textbf{Loss \& Representation Depth (Panel A)}: Cross-entropy loss descends smoothly from $5.60 \to 2.38$. Frozen linear probe cross-entropy evaluated across intermediate layers confirms monotonic refinement through depth ($\text{Block}_0: 3.9658 \to \text{Block}_1: 3.9630 \to \text{Block}_2: 3.9629$), proving that second-order rotation preserves signal fidelity.
    \item \textbf{Covariance Spectrum (Panel B)}: Minimum eigenvalues ($\lambda_L, \lambda_R$) remain non-negative ($\approx 10^{-14}$) under progressive ridge regularization, maintaining $0\%$ reconstruction residual.
    \item \textbf{Cautious Clip Dynamics (Panel C)}: Post-warmup clip activity stabilizes within a $2.5\% - 4.0\%$ band, preventing gradient surges from destabilizing shared residual paths.
    \item \textbf{Substrate Contribution (Panel D)}: Cross-branch parameters provide a sustained positive refresh contribution ($\Delta h_{\text{Ref}} = +2.3 \times 10^{-2}$).
\end{enumerate}

\subsection{Honest Taxonomy: Wins, Losses, and Trade-offs}

\begin{tcolorbox}[title=\textbf{RACASO Applicability Matrix}]
\textbf{WHEN RACASO IS OPTIMAL:}
\begin{itemize}
    \item Branching \& coupled multi-stream architectures where covariance goes rank-1.
    \item Loss landscapes with pronounced saddle points requiring negative curvature detection.
    \item Graphs with potential second-derivative operator hazards (protected by L5).
\end{itemize}
\textbf{WHEN STANDARD DIAGONAL OPTIMIZERS ARE PREFERABLE:}
\begin{itemize}
    \item Standard feedforward Transformers/ResNets with heavy normalization (AdamW).
    \item Memory-constrained environments where rotation matrices exceed budget (Lion).
    \item Pure convex classification problems where Sophia clipping slows early progress.
\end{itemize}
\end{tcolorbox}

\section{Usage \& Integration}

RACASO is implemented as a self-contained PyTorch optimizer:

\begin{verbatim}
import torch
from racaso import RACASO

model = MyModel().cuda()
params = list(model.parameters())

def forward_fn(params):
    return model(batch).loss  # Re-evaluable loss for Hutchinson HVP

# Hutchinson Mode (True Second-Order)
opt = RACASO(
    params,
    lr=3e-4,
    betas=(0.965, 0.99),
    rho=0.04,
    gamma=0.04,
    curvature_mode="hutchinson",
    forward_fn=forward_fn,
    hessian_freq=10,
    refresh_freq=10,
)

# Standard training loop
for batch in dataloader:
    loss = model(batch).loss
    loss.backward()
    opt.step()
    opt.zero_grad()
\end{verbatim}

For second-derivative-free training, configure \texttt{curvature\_mode="soap"}:
\begin{verbatim}
# SOAP Mode (Rotated Second Moment, No Forward Function Required)
opt = RACASO(params, lr=3e-4, curvature_mode="soap")
\end{verbatim}

\section{Conclusion}

RACASO establishes a principled framework for combining Kronecker-factored structural preconditioning with cautious second-order optimization. By evaluating curvature directly in the privileged eigenbasis and bounding updates through a five-layer safety architecture, RACASO resolves the numerical failure modes of earlier second-order methods. While standard diagonal optimizers remain optimal for simple feedforward networks, RACASO provides an effective, robust solution for complex branching architectures and coupled loss landscapes.

\section*{Acknowledgments \& Open-Source Lineage}

We express our sincere gratitude to Dr.~Ben Goertzel for his invaluable guidance, discussions, support, and arXiv endorsement.

We gratefully acknowledge the foundational open-source contributions and architectures that inspired RACASO:
\begin{itemize}
    \item \textbf{CASPR} (Cautious Adaptive Second-order Preconditioned Regularizer) for the foundational architectural intuition bridging second-order cautious bounds with structured preconditioning.
    \item \textbf{SOAP} (Vyas et al., 2024, Harvard/MIT/OpenAI) for Kronecker-factored structural preconditioning in the eigenbasis (MIT License).
    \item \textbf{Sophia} (Liu et al., 2023, Stanford) for per-element cautious clipping and Gauss-Newton-Bartlett estimation (MIT License).
    \item \textbf{Shampoo} (Gupta et al., 2018, Google Research) for tensor preconditioning principles (Apache 2.0).
    \item \textbf{AdaHessian} (Yao et al., 2020, UC Berkeley) for second-order adaptive learning rate scaling (Apache 2.0).
    \item \textbf{RAdam} (Liu et al., 2019) for variance confidence cold-start gating (Apache 2.0 / MIT).
    \item \textbf{Yogi} (Zaheer et al., 2018, Google Research) for additive variance bounds under bursty gradients (Apache 2.0 / MIT).
\end{itemize}

\bibliographystyle{plain}
\begin{thebibliography}{10}

\bibitem{gupta2018shampoo}
V.~Gupta, T.~Koren, and Y.~Singer.
\newblock Shampoo: Preconditioned stochastic tensor optimization.
\newblock {\em arXiv:1802.09568}, 2018.

\bibitem{hutchinson1989}
M.~F. Hutchinson.
\newblock A stochastic estimator of the trace of the influence matrix for laplacian smoothing splines.
\newblock {\em Communications in Statistics---Simulation and Computation}, 18(3):1059--1076, 1989.

\bibitem{kingma2014adam}
D.~P. Kingma and J.~Ba.
\newblock Adam: A method for stochastic optimization.
\newblock {\em arXiv:1412.6980}, 2014.

\bibitem{liu2023sophia}
H.~Liu, Z.~Li, D.~Hall, P.~Liang, and T.~Ma.
\newblock Sophia: A scalable stochastic second-order optimizer for language model pre-training.
\newblock {\em arXiv:2305.14342}, 2023.

\bibitem{liu2019radam}
L.~Liu, H.~Jiang, P.~He, W.~Chen, X.~Liu, J.~Gao, and J.~Han.
\newblock On the variance of the adaptive learning rate and beyond.
\newblock {\em arXiv:1908.03265}, 2019.

\bibitem{loshchilov2017adamw}
I.~Loshchilov and F.~Hutter.
\newblock Decoupled weight decay regularization.
\newblock {\em arXiv:1711.05101}, 2017.

\bibitem{martens2015kfac}
J.~Martens and R.~Grosse.
\newblock Optimizing neural networks with kronecker-factored approximate curvature.
\newblock {\em arXiv:1503.05671}, 2015.

\bibitem{pearlmutter1994fast}
B.~A. Pearlmutter.
\newblock Fast exact multiplication by the hessian.
\newblock {\em Neural Computation}, 6(1):147--160, 1994.

\bibitem{vyas2024soap}
N.~Vyas, D.~Morwani, R.~Zhao, G.~Kaplun, S.~Kakade, and B.~Barak.
\newblock Soap: Improving and stabilizing shampoo using adam.
\newblock {\em arXiv:2409.11321}, 2024.

\bibitem{yao2020adahessian}
Z.~Yao, A.~Gholami, S.~Shen, M.~Mustafa, K.~Keutzer, and M.~W. Mahoney.
\newblock Adahessian: An adaptive second order optimizer for machine learning.
\newblock {\em arXiv:2006.00719}, 2020.

\bibitem{zaheer2018yogi}
M.~Zaheer, S.~J. Reddi, D.~Sachan, S.~Kale, and S.~Kumar.
\newblock Adaptive methods for nonconvex optimization.
\newblock In {\em NeurIPS}, 2018.

\end{thebibliography}

\end{document}
